\documentclass[10pt]{article}
\usepackage[a4paper,margin=26mm]{geometry}
\usepackage[T1]{fontenc}
\usepackage{lmodern}
\usepackage{amsmath,amssymb}
\emergencystretch=2em
\title{A sparse penalty whose proximal map is not soft thresholding\\Disproof of Conjecture 00000008855}
\author{ziangni-sys}
\date{4 October 2026}
\begin{document}
\maketitle
\section*{Penalty and proximal definition}
The source asserts that proximal operators of sparse penalties are
always soft thresholding, without restricting the penalty to a convex
class. Consider the standard scalar nonzero-count penalty
\[
P(y)=\|y\|_0=
\begin{cases}0&y=0,\\1&y\ne0.\end{cases}
\]
This is a sparsity penalty: it charges one for a nonzero coordinate and
zero for a zero coordinate.
With unit proximal parameter, define the proximal graph by global
minimization of
\[
F_x(y)=\tfrac12(y-x)^2+P(y),\qquad
\operatorname{prox}_{P}(x)=\operatorname*{argmin}_{y\in\mathbb R}F_x(y).
\]
It may be a set in general, but the two inputs below have unique minimizers.

\section*{Two exact global minima}
At $x=2$, $F_2(2)=1$ whereas $F_2(0)=2$.
For every nonzero $y$, $F_2(y)=1+(y-2)^2/2\geq1$,
with equality if and only if $y=2$. Thus
$\operatorname{prox}_{P}(2)=\{2\}$.
At $x=1$, $F_1(0)=1/2$. For every nonzero $y$,
$F_1(y)=1+(y-1)^2/2\geq1>1/2$.
Consequently $\operatorname{prox}_{P}(1)=\{0\}$.
These comparisons check every real competitor; they do not rely on
an assumed hard-thresholding table.

\section*{No fixed soft threshold can realize this map}
For a threshold $\tau\geq0$, the standard scalar soft threshold is
\[
S_\tau(x)=
\begin{cases}
x-\tau&x>\tau,\\
x+\tau&x<-\tau,\\
0&|x|\leq\tau.
\end{cases}
\]
If $S_\tau$ were a proximal map for this penalty, uniqueness at $x=2$
would require $S_\tau(2)=2$.
When $\tau\geq2$ its value is zero; when $0\leq\tau<2$ its value
is $2-\tau$. The required equality therefore forces $\tau=0$.
But $S_0(1)=1$, contradicting the unique proximal value zero at $x=1$.
Thus no fixed nonnegative threshold can represent the proximal operator.
In particular the customary unit threshold gives $S_1(2)=1$,
which is not a minimizer.

This refutes the first clause of the original conjunction; the separate
splitting-rate clause is not needed for the disproof.
Soft thresholding for the absolute-value penalty is a different result
and does not imply the universal statement for sparse penalties.

\section*{Formal verification}
Lean defines the actual piecewise nonzero-count penalty, the real
quadratic objective, and the proximal relation as comparison with every
real competitor. It proves both exact proximal-set characterizations,
defines the piecewise soft-threshold function, and proves the impossibility
of a threshold valid for every input.
The final theorem \texttt{conjecture\_8855\_false} is stronger than a
failure at one prescribed threshold: even changing the fixed threshold
cannot fix this penalty.
\end{document}
